% Píše: backend.
\chapter{Hardware a~zapojení}
\label{chap:hardware}

Dům je sestavený z~modulů na nepájivém poli, které jsou připojené k~Raspberry Pi 3B+.
Raspberry Pi má čtyřjádrový 64bitový procesor Broadcom BCM2837B0 s~frekvencí 1,4~GHz, vestavěný
Bluetooth 4.2 s~podporou BLE a~rozšířenou 40pinovou lištu GPIO; napájí se zdrojem 5~V přes
konektor micro-USB a~výrobce uvádí proud 2,5~A~\cite{rpi3bplus}.

\section{Součástky}
\label{sec:soucastky}

Tabulka~\ref{tab:soucastky} shrnuje součástky konečné sestavy. Čísla GPIO jsou v~číslování
BCM (Broadcom), které používá i~program; fyzické piny lišty jsou v~tabulce~\ref{tab:zapojeni}.

\begin{table}[H]
\centering
\caption{Součástky domu}
\label{tab:soucastky}
\small
\begin{tabular}{|p{3.2cm}|p{3.9cm}|p{3.3cm}|p{2.5cm}|}
\hline
\textbf{Součástka} & \textbf{Úloha v~domě} & \textbf{Rozhraní} & \textbf{Napájení} \\
\hline
Raspberry Pi 3B+ & řídicí jednotka, periferie BLE & -- & 5~V, micro-USB \\
\hline
BMP280 (prodávaný jako BME280) & venkovní teplota a~tlak & I\textsuperscript{2}C, adresa
  \kod{0x76} & 3,3~V \\
\hline
TSL2561 & venkovní osvětlení & I\textsuperscript{2}C, adresa \kod{0x39} & 3,3~V \\
\hline
displej SSD1306, OLED 128~$\times$~64 bodů & stav domu & I\textsuperscript{2}C, adresa
  \kod{0x3C} & 3,3~V \\
\hline
HC-SR501 & čidlo pohybu (PIR) & vstup GPIO24 & 5~V, výstup 3,3~V \\
\hline
modul LED semafor & kontrolky topení, chlazení a~pohybu & výstupy GPIO17, GPIO27, GPIO22 &
  z~výstupů GPIO \\
\hline
RGB LED se společnou katodou & světlo domu & výstupy PWM GPIO16, GPIO20, GPIO21 & z~výstupů
  GPIO přes rezistory \\
\hline
\end{tabular}
\end{table}

Čidla měří skutečné hodnoty v~místnosti, kde sestava stojí, a~v~simulaci zastupují venkovní
prostředí (kapitola~\ref{chap:navrh}). Barvu a~jas RGB LED řídí pulzně šířková modulace (PWM).
Modul LED semafor má tři LED: červená ukazuje topení, zelená chlazení a~oranžová (prostřední)
pohyb.

\section{Zapojení}
\label{sec:zapojeni}

Tabulka~\ref{tab:zapojeni} uvádí zapojení podle fyzických pinů lišty Raspberry Pi. Piny, které
v~tabulce nejsou, zůstávají volné.

\begin{table}[H]
\centering
\caption{Zapojení podle fyzických pinů Raspberry Pi}
\label{tab:zapojeni}
\small
\begin{tabular}{|p{1.0cm}|p{3.2cm}|p{9.0cm}|}
\hline
\textbf{Pin} & \textbf{Funkce pinu} & \textbf{Připojený vývod} \\
\hline
1 & 3,3~V & VCC modulů BMP280, TSL2561 a~SSD1306 \\
\hline
2 & 5~V & VCC čidla HC-SR501 \\
\hline
3 & GPIO2 (SDA) & SDA modulů BMP280, TSL2561 a~SSD1306 \\
\hline
5 & GPIO3 (SCL) & SCL modulů BMP280, TSL2561 a~SSD1306 \\
\hline
6 & GND & GND čidla HC-SR501 \\
\hline
9 & GND & GND modulů BMP280, TSL2561 a~SSD1306 \\
\hline
11 & GPIO17 & R modulu LED semafor (topení) \\
\hline
13 & GPIO27 & G modulu LED semafor (chlazení) \\
\hline
15 & GPIO22 & Y modulu LED semafor (pohyb) \\
\hline
18 & GPIO24 & OUT čidla HC-SR501 \\
\hline
20 & GND & GND modulu LED semafor \\
\hline
34 & GND & společná katoda RGB LED (nejdelší vývod) \\
\hline
36 & GPIO16 & červená anoda RGB LED přes rezistor \\
\hline
38 & GPIO20 & zelená anoda RGB LED přes rezistor \\
\hline
40 & GPIO21 & modrá anoda RGB LED přes rezistor \\
\hline
\end{tabular}
\end{table}

Všechny tři moduly I\textsuperscript{2}C sdílejí sběrnici číslo~1 na pinech~3 (SDA) a~5 (SCL),
takže k~Raspberry Pi vedou jen dva datové vodiče. Adresy modulů se nepřekrývají. Raspberry Pi má
na GPIO2 a~GPIO3 pevné rezistory 1,8~k$\Omega$ připojené k~napájení~\cite{rpi-hardware}. Po
zapojení musí příkaz \kod{i2cdetect -y 1} na Raspberry Pi ukázat adresy \kod{39}, \kod{3c}
a~\kod{76}.

\section{Napájení}
\label{sec:napajeni}

Celá sestava se napájí z~Raspberry Pi, bez externího zdroje. Moduly I\textsuperscript{2}C mají
napájení 3,3~V z~pinu~1. Na žádný pin GPIO ani na sběrnici I\textsuperscript{2}C se nesmí dostat
5~V, jinak se Raspberry Pi poškodí. Jedinou součástkou napájenou z~5~V (pin~2) je čidlo pohybu
HC-SR501. Jeho výstup OUT vede do GPIO24, proto se před připojením změřil multimetrem: při
hlášení pohybu měl 3,3~V, takže ho jde připojit přímo. LED modulu semafor a~RGB LED svítí
přímo z~výstupů GPIO, každá přes svůj rezistor.

\section{Poznámky k~součástkám}
\label{sec:poznamky-soucastky}

\begin{description}
  \item[BMP280 prodávaný jako BME280.] Čidlo bylo koupené jako BME280, ale v~registru ID čipu
        \kod{0xD0} má hodnotu \kod{0x58}, kterou má BMP280~\cite{bmp280}. BME280 tam má
        \kod{0x60}~\cite{bme280}. BMP280 měří teplotu a~tlak, vlhkost ne. Program typ čidla
        pozná podle ID a~vlhkost hlásí jako neznámou (oddíl~\ref{sec:ovladace}). Čidlo má být
        co nejdál od Raspberry Pi, které hřeje a~zkreslovalo by měření.
  \item[TSL2561 nesmí mířit na RGB LED.] Čidlo osvětlení má vývod ADDR nezapojený, a~proto
        adresu \kod{0x39}. Kdyby svítilo na čidlo světlo domu, automatika by se rozkmitala:
        světlo se při šeru rozsvítí, čidlo naměří jeho světlo, šero skončí a~světlo zhasne.
  \item[Nastavení HC-SR501.] Propojka je v~poloze H, takže výstup zůstává v~logické 1, dokud se
        před čidlem někdo hýbe. Trimr doby je na minimu (výstup pak drží asi 3~s) a~trimr
        citlivosti uprostřed. Po zapnutí napájení se čidlo asi minutu ustaluje a~může hlásit
        pohyb samo. Čidlo má mířit tam, kde se při předvádění pohybují lidé.
  \item[Dvoubarevný displej.] Horních 16 řádků bodů displeje SSD1306 svítí žlutě, zbytek jinou
        barvou; rozložení textu se tomu přizpůsobuje (oddíl~\ref{sec:displej}).
  \item[LED semafor.] Modul má ochranné rezistory přímo na desce, takže se jeho vývody R, G a~Y
        zapojují přímo na GPIO. LED svítí v~logické 1.
  \item[RGB LED.] Holá LED se čtyřmi vývody má společnou katodu (nejdelší vývod) na GND a~každá
        barva má vlastní rezistor 220 až 330~$\Omega$; bez něj by LED přetížila výstup GPIO.
        Ovladač počítá se společnou katodou, LED se společnou anodou nepodporuje.
\end{description}

\section{Vývoj sestavy}
\label{sec:vyvoj-sestavy}

Původní návrh počítal s~jinými díly: s~čidlem BME280, které měří i~vlhkost, s~čidlem osvětlení
BH1750, s~kroužkem dvanácti adresovatelných LED WS2812B napájeným z~5~V přes převodník úrovní
a~s~hodinami reálného času DS3231. Konečná sestava je jiná a~program se jí přizpůsobil.
Místo BME280 přišel BMP280 bez čidla vlhkosti (viz výše) a~místo BH1750 je zapojené TSL2561,
pro které program dostal vlastní ovladač. Kroužek WS2812B nahradila obyčejná RGB LED řízená PWM,
která nepotřebuje převodník úrovní ani zvláštní knihovnu. Hodiny DS3231 v~sestavě nejsou:
Raspberry Pi bere čas ze sítě protokolem NTP (Network Time Protocol) a~bez internetu ho nastaví
tablet příkazem \kod{cas} (oddíl~\ref{sec:rozhrani-prikaz}). Kontrolky
topení, chlazení a~pohybu obstarává modul LED semafor. Program dnes počítá jen s~konečnou
sestavou a~volby pro jiné součástky nemá.
